% Algorithm 1 — integrity scoring.
% Source of truth: proctoring/rules.py :: compute_integrity_score
% Constants verified against the running system on 2026-09-27:
%   SEVERITY_WEIGHTS = {critical: 12.0, high: 6.0, medium: 3.0, low: 1.0}
%   REPEAT_GROWTH = 1.0, CO_OCCURRENCE_STEP = 0.25, CO_OCCURRENCE_CAP = 2.0
%   integrity_fail_below = 55 (core/config.py)
% Exercised by metrics/m1_integrity.py; the monotonicity property below is
% checked over 39 cumulative event sequences.
%
% Requires: \usepackage{algorithm} \usepackage{algpseudocode}

\begin{algorithm}[t]
\caption{Integrity scoring $c(E) \rightarrow \text{integrity}$}
\label{alg:integrity}
\begin{algorithmic}[1]
\Require Proctoring event set $E$; severity weights $w$; growth $\gamma\!=\!1.0$;
         co-occurrence step $\sigma\!=\!0.25$, cap $\kappa\!=\!2.0$
\Ensure  Integrity score in $[0,100]$ and a verdict
\Statex
\State $K \gets \{\,\text{kind}(e) : e \in E\,\}$ \Comment{distinct event kinds}
\State $d \gets 0$
\ForAll{$k \in K$}
    \State $E_k \gets \{\, e \in E : \text{kind}(e) = k \,\}$
    \State $s_k \gets \max_{e \in E_k} w[\text{severity}(e)]$
        \Comment{worst severity for this kind, not the sum}
    \State $d \gets d + s_k \cdot \bigl(1 + \gamma \ln |E_k|\bigr)$
        \Comment{repeats grow sub-linearly}
\EndFor
\State $d \gets d \cdot \min\bigl(\kappa,\; 1 + \sigma\,(|K| - 1)\bigr)$
    \Comment{distinct kinds co-occurring are worse than one kind repeating}
\State $\text{integrity} \gets \max\bigl(0, \min(100, \text{round}(100 - d))\bigr)$
\Statex
\If{$\text{integrity} < \tau_{\text{fail}}$} \Comment{$\tau_{\text{fail}} = 55$}
    \State \Return $(\text{integrity},\ \textsc{flag})$
\ElsIf{$\text{integrity} \geq 80$ \textbf{and} no $e \in E$ is \textsc{critical}}
    \State \Return $(\text{integrity},\ \textsc{clean})$
\Else
    \State \Return $(\text{integrity},\ \textsc{review})$
\EndIf
\end{algorithmic}
\end{algorithm}

% ---------------------------------------------------------------------------
% Prose for the body text. Three design decisions are worth stating explicitly,
% because each one replaced an earlier rule that produced a wrong ordering.
%
% (1) Worst-severity-per-kind, not a sum over events. Summing every event lets
%     a long run of one low-severity signal outweigh a single critical one. The
%     shipped rule takes the worst severity observed for each distinct kind and
%     then adds a sub-linear term for how often that kind recurred, so
%     persistence still counts but cannot dominate.
%
% (2) The logarithmic repeat term. With gamma = 1.0, the n-th occurrence of a
%     kind contributes s_k * ln(n) beyond the first. Six no-face events are
%     therefore materially worse than one, without six being six times one.
%
% (3) The co-occurrence multiplier. Independent evidence types appearing
%     together is stronger evidence than any one type repeating, so the
%     deduction is scaled by 1 + 0.25(|K| - 1), capped at 2.0. The cap exists
%     so that a candidate with many distinct minor signals cannot be scored
%     below one with a single critical signal.
%
% Resolved defect, for the limitations/validity discussion: an earlier version
% of this function halved the deduction for repeated events, which inverted the
% ordering. Six no-face events scored 79 (review) while a candidate with a
% phone, visible notes, a second voice and two off-screen glances scored 86
% (clean) — the more serious case scored higher. An intermediate draft using
% bounded saturation scored a fully-absent candidate 88 (clean), which was also
% wrong. The rule above is the third version, and metrics/m1_integrity.py now
% checks monotonicity over 39 cumulative event sequences: adding an event must
% never raise the score. The paper should describe this as a resolved defect
% with the re-audit attached, not as an open limitation.
